\chapter{闵可夫斯基场论与欧几里得场论及统计力学之间的联系}

本节简要回顾关于欧几里得路径积分以及从欧几里得关联函数中提取物理内容的技术问题。详细内容参见文献 [3,4] 或统计场论方面的教材，例如 [22,23]。

$D$ 维闵可夫斯基场论（$D-1$ 个空间维和一个时间维）通过解析延拓与 $D$ 维欧几里得场论相联系。在 Wick 转动下
\begin{equation}
\begin{aligned}
x_0 &\equiv t \to -i x_4 \equiv -i\tau , \\
p_0 &\equiv E \to i p_4 .
\end{aligned}
\label{eq:5.1}
\end{equation}

欧几里得约定为
\begin{equation}
\begin{aligned}
x_E^2 &= \sum_{i=1}^{4} x_i^2 = \mathbf{x}^2 - t^2 = -x_M^2 , \\
p_E^2 &= \sum_{i=1}^{4} p_i^2 = \mathbf{p}^2 - E^2 = -p_M^2 .
\end{aligned}
\label{eq:5.2}
\end{equation}

$D$ 维欧几里得场论与统计力学之间的联系，在采用费曼路径积分方法对场论进行量子化时最为清晰。这一联系构成了使用与研究统计力学系统相同的方法进行数值模拟的基础。二者之间的等价关系概要列于表~\ref{tab:1} 中。

\begin{table}[htbp]
\centering
\caption{欧几里得场论与经典统计力学之间的等价关系。}
\label{tab:1}
\begin{tabular}{ll}
\toprule
\textbf{欧几里得场论} & \textbf{经典统计力学} \\
\midrule
作用量 & 哈密顿量 \\
作用量单位 $\hbar$ & 能量单位 $\beta = 1/kT$ \\
\shortstack{振幅的费曼权重 \\ $e^{-S/\hbar} = e^{-\int L\,dt/\hbar}$} & 玻尔兹曼因子 $e^{-\beta H}$ \\
\shortstack{真空到真空振幅 \\ $\int D\phi\, e^{-S/\hbar}$} & \shortstack{配分函数 \\ $\sum_{\mathrm{conf.}} e^{-\beta H}$} \\
真空能量 & 自由能 \\
真空期望值 $\langle 0|O|0\rangle$ & 正则系综平均 $\langle O\rangle$ \\
时间编序乘积 & 普通乘积 \\
格林函数 $\langle 0|T[O_1 \ldots O_n]|0\rangle$ & 关联函数 $\langle O_1 \ldots O_n\rangle$ \\
质量 $M$ & 关联长度 $\xi = 1/M$ \\
质量间隙 & 关联函数的指数衰减 \\
无质量激发 & 自旋波 \\
正规化：截断 $\Lambda$ & 格距 $a$ \\
重整化：$\Lambda \to \infty$ & 连续极限 $a \to 0$ \\
真空的改变 & 相变 \\
\bottomrule
\end{tabular}
\end{table}

让我先从欧几里得理论的简要总结开始。考虑闵可夫斯基空间中的一般两点关联函数 $\Gamma_M(t,k) = \langle 0| \int d^3x\, e^{-ik\cdot x}\, T[F(x,t)F^\dagger(0,0)] |0\rangle$。海森堡算符 $F(x,t)$ 可以写成
\begin{equation}
F(x,t) = e^{iHt-ipx} F e^{-iHt+ipx}
\label{eq:5.3}
\end{equation}
于是
\begin{equation}
\Gamma_M(t,k) = \langle 0| \int d^3x\, e^{-ik\cdot x}\, e^{iHt-ipx} F e^{-iHt+ipx} F^\dagger |0\rangle
\label{eq:5.4}
\end{equation}
其中，在去掉时间编序乘积时，假设 $t > 0$。利用真空在时间与空间平移下的不变性，并为简洁起见假设 $F^\dagger = F$，则上式变为
\begin{equation}
\begin{aligned}
\Gamma_M(t,k) &= \langle 0| \int d^3x\, e^{-ik\cdot x}\, F e^{-iHt+ipx} F |0\rangle \\
&= \langle 0| \delta(p-k)\, F e^{-iHt} F |0\rangle .
\end{aligned}
\label{eq:5.5}
\end{equation}

由于时间依赖是显式的，我们可以利用 $t \to -i\tau$ 并保持 $H$ 不变，Wick 转动到欧几里得空间
\begin{equation}
\Gamma_E(t) = \langle 0| \delta(p-k)\, F e^{-H\tau} F |0\rangle .
\label{eq:5.6}
\end{equation}
为得到路径积分表示，我们将时间间隔划分为 $N$ 步 $\tau = Na$，并在每一步插入一组完备态 $|\phi\rangle\langle\phi|$。
\begin{equation}
\Gamma_E(t) = \delta(p-k) \sum_{\phi_1 \ldots \phi_N} \langle 0|F|\phi_1\rangle \langle \phi_1| e^{-Ha} |\phi_2\rangle \ldots \ldots \langle \phi_{N-1}| e^{-Ha} |\phi_N\rangle \langle \phi_N| F |0\rangle
\label{eq:5.7}
\end{equation}

关于通过引入完备态组，如何将量子力学振幅的算符形式换为对路径求和，这一讨论是标准的，因此不再重复。在这个「格点」上，传递矩阵 $T \equiv \exp(-Ha)$ 扮演时间平移算符的角色，由此可得连续哈密顿量为
\begin{equation}
H = \lim_{a\to 0} -\frac{1}{a} \log T .
\label{eq:5.8}
\end{equation}
为使 $H$ 成为自伴哈密顿量，要求 $T$ 是作用在具有正范数的态希尔伯特空间上的对称、有界且正的算符。这一点的必要条件就是 Osterwalder-Schrader 反射正性 [24]。LQCD 满足该条件已由 M.~L\"uscher 通过显式构造证明 [25]。因此，所需的物理闵可夫斯基理论中的关联函数可以通过对欧几里得对应关联函数作解析延拓而得到。

假定合理的欧几里得理论所需的条件已经确立，则在存在源 $J$ 时，真空到真空振幅（也称为生成泛函或配分函数）为
\begin{equation}
Z[J] = \int d\phi\, e^{-S_E + J\phi}
\label{eq:5.9}
\end{equation}
其中 $S_E$ 是欧几里得作用量。关联函数可以通过对 $J$ 微分 $\log Z[J]$ 的标准方式由此获得，例如 $\langle 0|\phi(x)\phi(0)|0\rangle = \frac{\delta}{\delta J(x)} \frac{\delta}{\delta J(0)} \log Z[J]\big|_{J=0}$。

最后，我总结欧几里得量与闵可夫斯基量之间的关系。基本工具是在式~\eqref{eq:5.7} 中反复插入一组完备的动量态。我采用态的相对论归一化
\begin{equation}
\langle k|p\rangle = 2E L^3\, \delta_{k_x p_x} \delta_{k_y p_y} \delta_{k_z p_z} \xrightarrow{L\to\infty} 2E(2\pi)^3 \delta^3(p-k) .
\label{eq:5.10}
\end{equation}
于是，对于与插值场 $F$ 耦合、质量为 $M_n$ 的稳定孤立态 $|n\rangle$ 谱，
\begin{equation}
\Gamma_E(\tau) = \sum_n \frac{\langle 0|F|n\rangle \langle n|F|0\rangle}{2M_n} e^{-M_n \tau} .
\label{eq:5.11}
\end{equation}
现在很容易得到这些 $M_n$ 与物理谱之间的联系。对欧几里得时间作傅里叶变换，随后转动到闵可夫斯基空间，给出
\begin{equation}
\int d\tau\, e^{ip_4 \tau}\, \frac{e^{-M_n |\tau|}}{2M_n}
= \frac{1}{2M_n (M_n - ip_4)} + \frac{1}{2M_n (M_n + ip_4)}
= \frac{1}{M_n^2 + p_4^2}
\xrightarrow{p_4 \to -iE} \frac{1}{M_n^2 - E^2} .
\label{eq:5.12}
\end{equation}

这些 $M_n$ 正是 $|n\rangle$ 传播子中极点的位置。这就是为什么可以简单地略过必要的解析延拓：从指数衰减测得的质量就是物理态的极点质量。当 $\tau \to \infty$ 时，由于指数压低，只有最低态有贡献，这提供了一种分离最轻质量态的方法。

这一简单联系对矩阵元同样成立。考虑 $\tau_2 \gg \tau_1 \gg 0$ 情形下的三点函数
\begin{equation}
\begin{aligned}
&\int d^3x\, d^3y\, e^{-ik\cdot x} e^{-iq\cdot y}\, \langle 0| T[F_f(x,\tau_2) O(y,\tau_1) F_i(0,0)] |0\rangle \\
&= \sum_a \sum_b \frac{\langle 0|F_f|a\rangle e^{-M_a (\tau_2 - \tau_1)}}{2M_a}\, \langle a|O|b\rangle\, \frac{\langle b|F_i|0\rangle e^{-M_b \tau_1}}{2M_b} .
\end{aligned}
\label{eq:5.13}
\end{equation}

$\langle a|O|b\rangle$ 左右两侧的两个因子，各相差一个衰变振幅，都与式~\eqref{eq:5.11} 中讨论的两点函数相同。因此，在 $\tau_2 - \tau_1$ 和 $\tau_1$ 很大时，它们提供了对态求和中的最低态的分离。这些都可以通过计算两点函数得到。因此，在欧几里得模拟中测得的 $\langle a|O|b\rangle$ 就是这些态之间所需的矩阵元，无需再解析延拓到闵可夫斯基空间。

一旦对态的求和不再仅仅是对稳定单粒子态求和，上述论证便不再成立。例如，考虑 $F$ 是 $\rho$ 介子插值场的情形。在这种情况下，还存在 $2\pi$、$4\pi$ 等中间态，其割线始于 $E = 2M_\pi$。$\rho$ 介子的物理特征（质量与宽度）可以用传播子中的一个复极点来表示，在闵可夫斯基空间中，该极点位于割线之下的第二黎曼叶上。这个物理极点不能再通过测量大 $\tau$ 处欧几里得关联函数指数衰减给出的「能量」来重现 [26]。它需要测量所有中间态与动量（所有 $\tau$）下的关联函数，然后对 $\tau$ 作傅里叶变换，再作解析延拓 $p_4 \to -iE$。

出于同样的原因，并不存在计算 $n$ 粒子散射振幅的简单方法。一般的两粒子散射振幅（四点关联函数）之所以是复数，正是因为中间态涉及对共振态以及所有动量取值的连续态的求和。原则上，这可以从欧几里得关联函数获得，但前提是以本质上无限的精度测量所有可能的 $\tau$ 下的关联函数，然后再作解析延拓。在实践中，这是毫无希望的可能。唯一的例外是阈上的振幅，它是实数。M.~L\"uscher 已经证明，阈上的两粒子散射振幅如何由有限体积中的能量移动来计算，即两粒子能量与两个孤立粒子能量之和的差 [27--30]。
